\documentclass[10pt,a4paper]{amsart}
\usepackage[margin=19mm]{geometry}
\usepackage{amsmath,amssymb,mathtools}
\usepackage[scr=rsfs,cal=euler]{mathalfa}
\usepackage{xfrac}
\usepackage[unicode,hidelinks,bookmarks=false]{hyperref}
\newcommand{\ie}{i.e.\ }
\newcommand{\cf}{cf.\ }
\newcommand{\resp}{respectively\ }
\newcommand{\Subsection}{\subsection}
\providecommand{\nobreakdash}{\nobreak\hbox{-}}
\setlength{\emergencystretch}{2em}
\multlinegap0pt
\usepackage{amssymb}
\usepackage{mathtools}
\usepackage{dsfont}
\usepackage[normalem]{ulem}
\usepackage{algorithm}\usepackage{algpseudocode}
\usepackage{csquotes}
\usepackage{tikz}
\usepackage[shortlabels]{enumitem}
\newtheorem{lemma}{Lemma}
\newtheorem{theorem}{Theorem}
\usepackage{cleveref}
\crefname{lemma}{Lemma}{Lemmas}
\crefname{theorem}{Theorem}{Theorems}
\title{Foundational Analysis Of The Solvability Complexity Index: The Weihrauch-SCI Intermediate Hierarchy}
\author{Christopher Sorg}
\date{Source: arXiv:2603.18955v3 (CC BY 4.0)}
\begin{document}
\maketitle

\noindent\emph{Source and licence.} Foundational Analysis Of The Solvability Complexity Index: The Weihrauch-SCI Intermediate Hierarchy. Version arXiv:2603.18955v3 (CC BY 4.0), distributed by arXiv under the Creative Commons Attribution 4.0 International licence, http://creativecommons.org/licenses/by/4.0/, as recorded in the arXiv licence metadata for this exact version and shown on the abstract page https://arxiv.org/abs/2603.18955v3. Full licence notice: \texttt{licenses/arxiv-2603.18955-CC-BY-4.0.txt}.\par\smallskip

\noindent\textit{Editorial scope.} Counted unit: the article's theorem thm:Xi-m-exact-height (source0.tex lines 2063--2068). The article's complete original proof of it is delivered with it (source0.tex lines 2070--2125, one \texttt{proof} environment of 56 lines). No mathematics is written, repaired or re-derived: the statement and the proof are the article's own lines, in the article's own notation, and no retained line is substituted or re-flowed.  Retained uncounted as the source-local context the proof needs: lines 1937--1946; lines 1981--1990; lines 2031--2033.  Nothing was omitted from any retained span, so the package carries no omission record.  No retained line contains a citation, so the wrapper carries no bibliography and no bibliography entry.  No retained line contains a figure environment, an image-inclusion command or drawing code, so no image is delivered and the package declares no asset dependency.  Editorial formatting only, in addition: an AMS \texttt{amsart} preamble that restates the article's own declarations and macros used by the retained lines, namely the environments \texttt{lemma}, \texttt{theorem} on separate counters, together with the standard packages those lines need.  The retained environments print in this document as Lemma 1, Theorem 1 in order of appearance, whereas the article prints the same items with the numbering of its own full document.  The complete native source of the article as distributed in the arXiv e-print for this exact version is bundled unchanged as \texttt{sources/196741/source0.tex}.
\begin{lemma}[Height \(k\) coordinate towers yield \(\Delta^0_{k+1}\) sets]\label{lem:height-k-delta}
Let
\[
\Theta:\Omega_1\to\{0,1\}
\]
be computed by a type-G tower of height \(k\) over the coordinate oracle \(\Lambda_{\mathrm{mat}}\). Then
\[
\Theta^{-1}(\{1\})\in\Delta^0_{k+1}(\Omega_1).
\]
\end{lemma}

\begin{lemma}[\(\Sigma^0_m\)-completeness]\label{lem:Xi-m-complete}
For every \(m\in \mathbb{N}\),
\[
\Xi_m^{-1}(\{1\})
\]
is \(\Sigma^0_m\)-complete in the Cantor space \(\Omega_1\). In particular,
\[
\Xi_m^{-1}(\{1\})\notin\Delta^0_m(\Omega_1).
\]
\end{lemma}

\begin{lemma}[Finite Boolean combinations preserve tower height]\label{lem:boolean-combinations-same-height}
Let \(\Theta_1,\ldots,\Theta_N:\Omega_1\to\{0,1\}\) be decision maps each computed by a type-G tower of height at most \(k\). Then every Boolean combination of \(\Theta_1,\ldots,\Theta_N\) is computed by a type-G tower of height at most \(k\).
\end{lemma}

\begin{theorem}[Exact finite type-G heights of the matrix source problems]\label{thm:Xi-m-exact-height}
For every \(m\in \mathbb{N}\),
\[
\mathrm{SCI}_G(\mathcal P_m^{\mathrm{mat}})=m.
\]
\end{theorem}

\begin{proof}
We prove the upper and lower bounds separately.

\smallskip
\noindent\textbf{Upper bound:}
For a finite tuple \(s=(n_1,\ldots,n_r)\), \(0\le r\le m\), define the tail predicate
\[
\Xi_{m,r}(A;s)
\]
as follows: If \(r=m\), put
\[
\Xi_{m,m}(A;n_1,\ldots,n_m) := A(\beta_m(n_1,\ldots,n_m)).
\]
If \(r<m\), define recursively
\[
\Xi_{m,r}(A;s)=1
\]
if and only if
\[
(Q_{r+1} n_{r+1})\cdots(Q_m n_m) \bigl[A(\beta_m(n_1,\ldots,n_m))=1\bigr].
\]
Thus \(\Xi_m(A)=\Xi_{m,0}(A;\varnothing)\).

We prove by backward induction on \(r\) that \(\Xi_{m,r}(\,\cdot\,;s)\) is computable by a type-G tower of height at most \(m-r\), uniformly in the finite parameter \(s\).

For \(r=m\), the map
\[
A\mapsto A(\beta_m(n_1,\ldots,n_m))
\]
is a single coordinate query, hence a height-\(0\) general algorithm. Assume now \(r<m\). If \(Q_{r+1}=\exists\), then
\[
\Xi_{m,r}(A;s) = \lim_{N\to\infty} \max_{1\le n\le N}\Xi_{m,r+1}(A;s,n).
\]
If \(Q_{r+1}=\forall\), then
\[
\Xi_{m,r}(A;s) = \lim_{N\to\infty} \min_{1\le n\le N}\Xi_{m,r+1}(A;s,n).
\]
For each fixed \(N\), the finite maximum or minimum is a finite Boolean combination of maps computable by towers of height at most \(m-r-1\). By \cref{lem:boolean-combinations-same-height}, it is computable by a tower of height at most \(m-r-1\). The additional limit over \(N\) gives height at most \(m-r\). Taking \(r=0\) gives a height-\(m\) tower for \(\Xi_m\).

\smallskip
\noindent\textbf{Lower bound:}
Suppose, toward a contradiction, that
\[
\mathrm{SCI}_G(\mathcal P_m^{\mathrm{mat}})\le m-1.
\]
Then \(\Xi_m\) is computed by a type-G tower of height at most \(m-1\). By \cref{lem:height-k-delta},
\[
\Xi_m^{-1}(\{1\})\in\Delta^0_m(\Omega_1).
\]
This contradicts \cref{lem:Xi-m-complete}, which says that \(\Xi_m^{-1}(\{1\})\) is \(\Sigma^0_m\)-complete and therefore not in \(\Delta^0_m\).
Thus
\[
\mathrm{SCI}_G(\mathcal P_m^{\mathrm{mat}})\ge m.
\]
Combining both inequalities proves the theorem.
\end{proof}

\end{document}
